% capitulos/conclusao.tex

Esta primeira etapa do projeto LibrasVision, desenvolvida no âmbito do TCC1, dedicou-se ao levantamento bibliográfico sistemático, à fundamentação teórica e ao planejamento metodológico do sistema de reconhecimento de sinais dinâmicos da Língua Brasileira de Sinais (Libras) para ambientes educacionais. O trabalho realizado nesta fase constitui a base conceitual e técnica sobre a qual o desenvolvimento, a implementação e a avaliação do sistema serão conduzidos na etapa subsequente (TCC2).

A revisão da literatura confirmou que o campo do reconhecimento automático de língua de sinais encontra-se em acelerada evolução, com resultados expressivos obtidos por arquiteturas baseadas em MediaPipe Holistic combinado a redes LSTM, que alcançaram acurácia de até 97\% em diferentes línguas de sinais \cite{anithadevi:25,varshini:25}. Esse panorama valida a escolha arquitetural do LibrasVision e demonstra que a abordagem planejada está alinhada com o estado da arte. Ao mesmo tempo, as revisões sistemáticas consultadas \cite{costa-slr:25,silva-ia:25} evidenciaram que a maioria dos sistemas disponíveis foi desenvolvida para línguas como a ASL (\textit{American Sign Language}) ou a ISL (\textit{Indian Sign Language}), com escassez de soluções voltadas especificamente à Libras em contexto universitário brasileiro, lacuna que o LibrasVision se propõe a preencher. A aceleração das pesquisas em acessibilidade digital para a comunidade surda, documentada por \citeonline{sousa:25}, reforça a relevância social e científica dessa contribuição.

O planejamento metodológico definido neste trabalho estabelece um pipeline estruturado em etapas bem delimitadas: definição do escopo com vocabulário acadêmico representativo; coleta de dados com voluntários em ambiente controlado; construção de um \textit{dataset} próprio com 150 vídeos em três classes; implementação do pipeline com OpenCV, MediaPipe Holistic e redes LSTM em TensorFlow/Keras \cite{kamble:25}; e avaliação do desempenho por métricas de acurácia, latência, F1-score e análise de robustez em condições variadas de iluminação, múltiplos usuários e oclusão parcial \cite{attili:24}. A definição de um protocolo rigoroso de avaliação comparativa com sistemas da literatura, incluindo o quadro comparativo planejado na Seção 3.6, permitirá situar os resultados do LibrasVision em relação ao que foi reportado nos principais trabalhos revisados.

Um diferencial importante da abordagem planejada é o compromisso com hardware convencional. Enquanto parte dos sistemas avaliados na literatura exige infraestrutura computacional especializada, como o uso de GPUs dedicadas \cite{fanucchi:24}, o LibrasVision será projetado para operar com uma webcam padrão em um computador convencional, viabilidade já demonstrada por \citeonline{attili:24} e \citeonline{kamble:25} em contextos similares. Essa característica é essencial para que a solução seja aplicável em instituições de ensino com recursos computacionais limitados, tornando a tecnologia acessível ao cenário universitário brasileiro onde a escassez de intérpretes é mais crítica \cite{dockens:18}.

O mapeamento antecipado das limitações é essencial para delimitar o escopo do projeto e orientar trabalhos futuros. O \textit{dataset} a ser construído abrangerá apenas sinais relacionados ao contexto acadêmico, o que restringirá a aplicabilidade do sistema a esse domínio específico. Além disso, a coleta com 4 voluntários em ambiente controlado poderá introduzir viés de representatividade: usuários com características físicas, velocidades de execução ou estilos de sinalização distintos dos presentes no \textit{dataset} poderão obter taxas de reconhecimento inferiores às projetadas, limitação documentada de forma análoga em \cite{abdulhussein:20}. A severidade desse risco é evidenciada por \citeonline{perdigoto:24}, que observaram queda de acurácia de 65,51\% para 6,97\% ao comparar imagens padronizadas com capturas de usuários leigos, reforçando a importância da diversidade na coleta. Adicionalmente, iluminação inadequada, oclusão parcial das mãos ou ângulos desfavoráveis da câmera poderão degradar a extração de \textit{landmarks} pelo MediaPipe, comprometendo a classificação posterior \cite{jain:24}. Essas limitações não invalidam a investigação proposta, mas delimitam o alcance das conclusões que poderão ser extraídas no TCC2.

Como próximos passos, o TCC2 deverá conduzir a implementação efetiva do pipeline, a coleta dos dados com os voluntários, o treinamento e a validação do modelo LSTM, e a avaliação experimental do sistema conforme o protocolo planejado neste trabalho. Os resultados obtidos serão comparados com os sistemas da literatura identificados na revisão bibliográfica, possibilitando verificar se o LibrasVision atinge a meta de acurácia superior a 85\% e latência inferior a 100 ms em condições próximas ao uso real em laboratórios e salas de aula universitárias \cite{costa-i3d:25}. Caso os objetivos sejam alcançados, o trabalho contribuirá concretamente para a redução das barreiras de comunicação enfrentadas por estudantes surdos no ensino superior brasileiro \cite{dockens:18}, demonstrando a viabilidade de tecnologias assistivas de baixo custo baseadas em visão computacional e aprendizado profundo.
